% LSPV-4 REPORT::OUTQUAD64::2026-10-13T20:03:06+00:00
% !TEX program = xelatex
\documentclass[11pt]{article}
\usepackage[T1]{fontenc}
\usepackage{lmodern,amsmath,amssymb,amsthm,booktabs,xurl,hyperref}
\usepackage[margin=0.88in]{geometry}
\hypersetup{hidelinks}
\setlength{\emergencystretch}{3em}
\newtheorem{theorem}{Theorem}
\newtheorem{lemma}[theorem]{Lemma}
\newtheorem{corollary}[theorem]{Corollary}
\newcommand{\F}{\mathbb F_2}
\newcommand{\SAT}{\operatorname{SAT}}
\title{OUTQUAD64: A Quadratic-Output Identity\\
\large Raw-coordinate inversion of the hidden map $u^3+u^8$}
\author{UNIVERSAL\_CORE\_BRIDGE / LSPV-4}
\date{14 October 2026, KST}
\begin{document}\maketitle
\paragraph{Status.}
\textbf{CANDIDATE with a complete symbolic proof; independent research review pending.}
A deterministic polynomial-time procedure, given only raw quadratic Boolean
coefficients, solves every target of an affine-hidden $u^3+u^8$ map using an
explicitly acquired affine enclosure of at most 64 points. On arbitrary
quadratic maps the same procedure is sound and partial. The statement is a
constructive upper bound, not a separation. \textbf{P versus NP remains OPEN.}

\section{The precise repair to BILINCUBE1}
BILINCUBE1 recovered identities affine in both input and output, obtaining an
all-target four-point bound for hidden $u^3+\lambda u$. Its nearby control
$u^3+u^8$ over $\mathbb F_{256}$ had six roots at zero, so that bound could not
apply unchanged. The present repair adds only \emph{pure output-quadratic}
terms, not products involving two unknown input bits.

Fix a binary field $K$ of degree $d\ge1$ and a coordinate identification
$K\cong\F^d$. Let $A,C$ be invertible binary-linear maps and let
$\alpha,\beta\in K$. Consider
\[
 P(z)=C\big((Az+\alpha)^3+(Az+\alpha)^8\big)+\beta.       \tag{1}
\]
The solver receives $d$, the complete squarefree quadratic coefficient table
of $P$, and $b\in\F^d$. It is \emph{not} given $K$'s modulus, $A,C,\alpha,\beta$,
or a source preimage. Arithmetic in the proof's hidden coordinates is field
arithmetic; arithmetic in the solver is Boolean coefficient matching and
binary linear algebra.

The raw encoding orders coefficients as constant, $d$ linear terms and
lexicographic pairs $i<j$, with each coefficient a $d$-bit output vector.
The implementation serializes these vectors as nonnegative JSON integers and
validates the explicit coefficient count before dimension-indexed allocation.
A dense fixed-width binary table has $\Theta(d^3)$ bits; the JSON version is
polynomially related. No succinct circuit expansion is free.

\section{An exact identity valid in every binary field}
\begin{lemma}[Quadratic-output identity]\label{lem:id}
Whenever $v=u^3+u^8$ in a field of characteristic two,
\[
 u^{64}+u^4+vu^{16}+v^2u^8+vu+v^3+v^8=0.               \tag{2}
\]
For fixed $v=b$, its left side is an affine binary-linear function of $u$,
and the corresponding polynomial has at most 64 roots.
\end{lemma}
\begin{proof}
Substitution gives
\begin{align*}
vu^{16}&=u^{19}+u^{24}, &v^2u^8&=u^{14}+u^{24},\\
vu&=u^4+u^9, &v^3&=u^9+u^{14}+u^{19}+u^{24},\\
v^8&=u^{24}+u^{64}.
\end{align*}
Each monomial cancels in (2), proving an identity in the ordinary polynomial
ring $\F[u]$, not merely at finitely many tested field points.
For fixed $b$, define
\[
 \mathcal L_b(u)=u^{64}+bu^{16}+b^2u^8+u^4+bu.
\]
The desired equation is $\mathcal L_b(u)=b^3+b^8$. Frobenius powers are
binary-linear, so its solutions form an affine space or the empty set.
As an ordinary polynomial $\mathcal L_b(u)+b^3+b^8$ is monic of degree 64.
The elementary polynomial root bound gives at most 64 roots, for all $d$,
including fields of fewer than 64 elements.
\end{proof}

The new term $v^3$ has Boolean degree two, whereas $v^8$ is binary-linear.
The mixed terms are still bilinear in the input and output coordinates.
Thus independent invertible affine changes in input and output keep every
coordinate of (2) in the following template:
\[
 R(z,y)=c+\sum_i a_i z_i+\sum_j b_j y_j+
 \sum_{i,j}M_{ij}z_i y_j+\sum_{j<k}D_{jk}y_jy_k.         \tag{3}
\]
No field representation needs to be recovered to search this template.

\section{Acquiring the identities and solving every target}
Let $\mathcal I(P)$ be the vector space of all coefficient arrays in (3)
satisfying $R(z,P(z))\equiv0$ in the Boolean ring
$\F[z_1,\ldots,z_d]/(z_i^2-z_i)$. The template has
\[
 c_d=(d+1)^2+\binom d2                                  \tag{4}
\]
unknowns. After substitution, Boolean degree is at most four, so matching
coefficients gives at most
\[
 r_d=\sum_{j=0}^{\min(4,d)}\binom dj                    \tag{5}
\]
linear equations. Compute a basis of the entire nullspace, then recheck each
basis element against every stored coefficient row. Absent rows are zero.
This uses the explicit ANF, not random evaluations or a sampled graph table.

For target $b$, substitute $y=b$ in every basis identity and solve the resulting
affine-linear system in $z$. Call its solution set $E_b$.
\begin{theorem}[Raw-coordinate 64-point selector]\label{thm:main}
One deterministic uniform algorithm, given only the raw coefficient table of
any map (1) and a target $b$, outputs the complete set $P^{-1}(b)$ in polynomial
bit time and polynomial workspace. It acquires $E_b$ by (3)--(5), and evaluates
at most 64 candidate inputs. No recognition algorithm for membership in the
hidden-map orbit is assumed.
\end{theorem}
\begin{proof}
Every graph identity vanishes at a real solution, hence $P^{-1}(b)\subseteq E_b$.
Write the hidden variables as $u=Az+\alpha$ and
$v=C^{-1}(y+\beta)$. Every binary coordinate of (2), expressed in $(z,y)$,
belongs to (3) and vanishes identically on the graph of $P$. It therefore
belongs to the computed full nullspace, whether or not the solver identifies
it as a field-derived relation. After fixing $y=b$, the equations defining
$E_b$ imply all coordinates of
$\mathcal L_{C^{-1}(b+\beta)}(Az+\alpha)=
(C^{-1}(b+\beta))^3+(C^{-1}(b+\beta))^8$.
Lemma~\ref{lem:id} and the bijectivity of $z\mapsto Az+\alpha$ imply
$|E_b|\le64$. Binary Gaussian elimination either reports inconsistency or
returns a particular point and a basis of dimension at most six. Enumerate
that affine set and retain precisely points satisfying the supplied $P(z)=b$.
This returns all and only the actual preimages, including correct empty
answers for targets outside the image.
\end{proof}

\paragraph{Sound partial behavior outside the promise.}
For any syntactically valid quadratic map, identities still imply
$P^{-1}(b)\subseteq E_b$. If $E_b$ is inconsistent or has dimension at most six,
the algorithm returns an exact answer after checking its entire enclosure.
If its dimension is larger, it returns \texttt{UNKNOWN}, not UNSAT and not an
assertion that the map lies outside the hidden orbit. Malformed maps and
out-of-range targets return \texttt{INVALID}. This is not a claimed total
decider of the language consisting of recognized hidden-map encodings.

\paragraph{Costs.}
Matrix dimensions are $O(d^4)$ by $O(d^2)$, with $O(d^6)$ dense matrix bits.
Schoolbook elimination uses $O(d^8)$ elementary Boolean coefficient operations;
the direct coefficient products also have polynomial construction cost.
Address arithmetic adds ordinary indexing overhead, not free large-word
operations. A conservative multitape realization by explicit scans is bounded
by $O(d^{12})$ bit time and polynomial workspace. The target substitution,
linear solve, and at most 64 raw-map evaluations are polynomial. Neither these
loose upper bounds nor Python timings are small-space certificates. They do
not establish any forbidden joint time-space point.

\section{Coupled Tseitin witnesses and the exact scope}
For a supplied simple graph $X$, charge vector $c$, and the deterministic
ordered cotree $C_X$ of size $d$, consider
\[
 [B_Xx=c]\ \land\ [P(\pi_{C_X}x)=b].                   \tag{6}
\]
The parser validates endpoints, edge uniqueness, explicit charges and cotree
dimension. Isolated vertices remain in the graph. Inconsistent component
charges give UNSAT; otherwise the forest lift produces
\[
 B_Xa=c,\quad B_XL=0,\quad\pi_{C_X}a=0,
 \quad\pi_{C_X}L=I,\qquad x=a+Lz.                       \tag{7}
\]
For every root returned by Theorem~\ref{thm:main}, reconstruct and verify the
complete edge witness. Conversely every edge witness projects to one returned
root. Forest discovery, charge adjustment, reconstruction, source access,
output and verification are charged. The implementation returns every witness;
a decision or single-witness consumer may stop after one verified root.

For arbitrary edge support $R$, retain the existing exact identity
\[
 \rho=\dim\pi_R\ker B_X=|R|+\operatorname{cc}(X)
                        -\operatorname{cc}(X-R),          \tag{8}
\]
where deletion keeps the whole vertex set, including isolates. Its
$2^{\rho+h}\operatorname{poly}(N)$ enumeration bound with $h$ external bits
is an upper bound, not an unavoidable-work statement. The present procedure
avoids that enumeration on the specific cotree residual (6).

An arbitrary source conjunct $F(w)$ is not solved by extending a supplied $w$.
Different easy conditions on a shared cotree may be incompatible. Finding a
source assignment and finding its extension remain different obligations.
Exact-(2,2) normalization may be applied only with its own charged encoding
and decoding interface; this report does not give a recognizer for arbitrary
anonymized normalized CNFs or transfer structured minimality to a new source.

\section{Reproducible experiments and falsification attempts}
The full record has 57,446 assertions in 15 independent invocations,
zero failures, total recorded time 25.823 seconds and longest invocation
17.436 seconds. Each audit invocation used a 25-second cooperative cap and
a 30-second external cap; no timeout occurred. These are \textbf{COMPUTATIONAL}
results, not asymptotic evidence. The symbolic theorem rests on its proof.

The small suite exhausts all 71 irreducible binary moduli of degrees one
through eight, using one seeded hidden affine map per modulus and all 10,824
targets. It compares complete root sets with brute force. It does not exhaust
all affine conjugations. Three larger maps at $d=16,24,32$ use twelve targets
each, with a different oracle: field-polynomial gcd with $Y^{2^d}-Y$ followed
by absolute-trace factor splitting. The raw solver never receives that oracle
or its secrets. A separate GF256 smoke suite tests 512 targets of canonical
and affine-hidden maps. A new-process isolation test contains only the
standalone solver and raw input JSON; no field or fixture module is present.

\begin{center}
\begin{tabular}{rrrrr}\toprule
$d$ & Old identities & New identities & Preprocessing (s) & Tested enclosure dimensions\\\midrule
16 & 1 & 17 & 0.232 & $0,2,4$, or empty\\
24 & 1 & 25 & 2.010 & $0,2,4$, or empty\\
32 & 1 & 33 & 16.884 & $0,2,4$, or empty\\\bottomrule
\end{tabular}
\end{center}
At $d=32$ the prior identities leave dimension 32 at the translated zero
target and dimension 31 at the other tested targets. The new system leaves
at most four dimensions there. The theorem guarantees at most six, not an
empirically inferred four-dimensional bound.

\paragraph{The exact preceding counterexample is resolved.}
For $K=\F[T]/(T^8+T^4+T^3+T+1)$ and $P(u)=u^3+u^8$, the old space has five
identities and the new space thirteen. At target zero, the new affine
enclosure has sixteen points. Evaluation returns exactly
\[
 \{0,1,12,80,176,237\},                                  \tag{9}
\]
using coefficient-bit integer encoding. Across all 256 targets, the output is
171 SAT and 85 UNSAT. Eighty of the latter are rejected directly by inconsistent
identities; five have a spurious singleton eliminated by verification.

\paragraph{A useful limitation, not hardness.}
Two independent copies of the same GF256 map have 36 roots at target $(0,0)$
and an eight-dimensional identity enclosure. The global six-dimensional cap
returns UNKNOWN. Recognizing the independent factors permits separate solves;
this easy product is not evidence of a machine lower bound. Constant,
projection, affine/parity, independent AND and generic quadratic controls
likewise receive only sound exact answers or explicit UNKNOWN.

The graph audit reconstructs 59 complete witnesses, checks 37 kernel columns,
four particular solutions and 48 projected ranks including isolates. It tests
four inconsistent charged isolates, sixteen malformed coupled records and a
shared-target incompatibility. Raw-map tests reject eight malformed maps and
four invalid targets. Full seeds, individual rows, timestamps and SHA-256
hashes are retained in the reproduction package and audit summary.

\section{Literature and theorem-transfer audit}
Patarin's input-output linearization and Diene et al.'s study of its relation
space are the methodological context \cite{patarin,diene}. The new task-specific
step here is the explicit identity (2) and its all-target bound in the expanded
template (3). No novelty claim about a general cryptanalytic method is made.
Gorla, Mueller and Petit study broader upper bounds for HFE polynomial systems
\cite{hfe}; those results are context, not an unverified premise of this proof.
The primary publisher/author abstracts and metadata were consulted in this
run; subscription full texts were not claimed read. All load-bearing algebra
in this report is given explicitly.

De Colnet and Mengel's Tseitin result concerns DNNF representation size and
regular-Resolution refutations \cite{tseitin}. A successful witness computation
need not construct either object. It cannot be turned into an arbitrary-machine
lower bound by counting all semantic states. The present fixed program has a
polynomial execution on the stated promise family even with hidden coordinates
and polynomial workspace. Its program description does not grow with that
workspace.

The upper bound survives stronger oracle access and supplies no
nonrelativizing or nonalgebrizing lower-bound mechanism. Natural Proofs is not
invoked. For the post-quantum ledger, input-output relation spaces give an
\textbf{EXACT OBJECT/THEOREM TRANSFER} at the constructive level. Treating
large rank, nonlinear support, collision counts or cryptographic conjectures
as unconditional runtime lower bounds remains a \textbf{FAILED TRANSFER}.
No ETH, GRH, cryptographic or average-case MQ hypothesis is used.

\section{Corpus, integration and publication record}
\paragraph{CORPUS REVIEWED.}
The live AGENTS.md, CODEX.md, research skill, redirect directive, root main.tex
and main4/main5 driver passages were fetched. The current recursive tree was
surveyed, though its displayed response was truncated; its root was
\path{92e5412bf9b5a4303d9a38d33c4492552523901f}, matching the previous read
before new publication. Relevant prior canonical MAIN, peer, task and
EARLY22/SINGLE22 bodies already read in this conversation were reused within
that unchanged snapshot rather than claimed freshly reread in full.
The exact local BILINCUBE1 report, solver, fixture, graph and boundary audit
bodies were inspected. A repository search for BILINCUBE1 returned no result,
consistent with its recorded local-only publication. No full-checkout preflight
or repository-wide validation is claimed.

\paragraph{DUPLICATION / CONTRADICTION AUDIT.}
BILINCUBE1 already provides raw coefficient matching and the four-point cube
bound. Its $u^3+u^8$ control is the precise input to this run, not a new
counterexample. OUTQUAD64 adds pure output-quadratic columns, identity (2),
and the proved 64-point all-target bound. The bounded search did not establish
literature-wide novelty. No prior compiler or restricted proof theorem is
contradicted.

\paragraph{EARLY22 / SINGLE22 acknowledgment.}
\textbf{ALREADY INTEGRATED, not extended beyond scope.} Refer to BILINCUBE1's
exact retained ledger and the earlier IMAGE22/ARC22 reports. LBO3SAT-0031/0032
and PROOF-0034 retain their signed exact-(2,2), minimality and Resolution
hypotheses; the clause-monotone cubic-pair map is separate. Historical variable
counts are not encoded bit lengths. The affine continuation limitations are
not reopened. The manuscript location of the new interface is Sections 3--4
here. Its next check is parameter-uniform extension of the new identity, not
an arbitrary-decider-to-proof extraction.

\paragraph{WHAT CHANGED RELATIVE TO ROOT MAIN.}
The live root was the markerless string \texttt{PLACEHOLDER\_FOR\_FULL\_MAIN},
blob \path{1cabdc5356780d0397cb47b865c525610f42c089}. It supplies no ownership
region to reconcile. No replacement or reconstruction is authorized for this
specialist. The scientific change relative to the last intact record is a new
upper-bound escape for hidden $u^3+u^8$; the missing unrestricted lower-bound
implication is unchanged.

\paragraph{ADOPTED / REJECTED / RETIRED PRIOR COMPONENTS.}
Adopted: exact raw-ANF identities, charged binary elimination, forest lifting
and off-promise UNKNOWN. Rejected: treating hidden field recovery as necessary
for this subtype. The proposal that the eight-power perturbation alone defeats
all polynomial selectors is \textbf{FALSIFIED/RETIRED} on the stated affine
orbit. General quadratic SAT, conjoined source search and common-input UWI
amplification remain \textbf{OPEN}.

\paragraph{Publication boundary.}
Exact code, raw outputs, report and prior report bodies were preserved locally.
GitHub write actions became available in this run. Canonical write/readback
outcomes are recorded separately in the outbox manifest, not inferred from
local success. Root reconciliation must wait for a legitimate LSPV-5 restoration
of unique ownership regions. Existing schedules and other owners are unchanged.

\paragraph{Canonical artifact paths.}
The standalone raw solver is
\path{experiments/review/outquad64_solver.py}; its independently coded finite
GF256 smoke audit is \path{experiments/review/outquad64_smoke_audit.py}.
The aggregate record is
\path{results/raw/review/2026-10-14_0503_outquad64_audit_summary.json}.
The full test harness, all rows and exact prior reports remain in the
reproducibility package. A separate implementation in the same run is not
independent research review.

\section{One next target}
Derive and test the dependence on a \emph{growing} exponent parameter $t$ for
hidden $u^3+u^{2^t}$, charging identity-template dimension and candidate count
as functions of both $d$ and $t$. Fixed $t=3$ is now an upper-bound escape;
raising a constant or observing failure of this particular solver does not
supply the required all-program lower bound. P versus NP remains OPEN.

\begin{thebibliography}{9}
\bibitem{patarin} J. Patarin. Cryptanalysis of the Matsumoto and Imai Public
Key Scheme of Eurocrypt'88. CRYPTO 1995, LNCS 963, 248--261.
\url{https://doi.org/10.1007/3-540-44750-4_20}.
\bibitem{diene} A. Diene, J. Ding, J. E. Gower, T. J. Hodges and Z. Yin.
Dimension of the Linearization Equations of the Matsumoto-Imai Cryptosystems.
WCC 2005 revised papers, LNCS 3969, 242--251, 2006.
\url{https://doi.org/10.1007/11779360_20}.
\bibitem{hfe} E. Gorla, D. Mueller and C. Petit. Stronger bounds on the cost
of computing Groebner bases for HFE systems. Journal of Symbolic Computation
109, 386--398, 2022. \url{https://arxiv.org/abs/2011.01050}.
\bibitem{tseitin} A. de Colnet and S. Mengel. Characterizing Tseitin-Formulas
with Short Regular Resolution Refutations. JAIR 76, 265--286, 2023.
\url{https://arxiv.org/abs/2103.09609}.
\end{thebibliography}
\end{document}
% END_OF_VERIFIED_REPORT::LSPV-4::OUTQUAD64
